\section{Flavor boundary and representation algebra}
\label{sec:flavour}

The electroweak mass-ratio problem is separated from the generation problem.  The working boundary condition is that generations are not assumed to arise exactly from compactification topology.  A separate SO(32)-flavor interpretation may organize fermion content and endpoints, and it may or may not be compatible with a compactification geometry.

Representation-theoretic unification remains useful because it constrains how Standard Model fermions can sit inside larger groups.  The algebraic GUT literature gives clean embeddings through SU(5), Spin(10), Pati--Salam, and E6 representation theory \cite{BaezHuerta2010,Britto2021}.  These embeddings can guide the flavor and endpoint sector without determining the DeVries W/Z ratio by themselves.

A future compactification model must satisfy two independent constraints.  First, it must reproduce the electroweak projective mass datum.  Second, it must accommodate the observed chiral matter and flavor structure.  The project does not identify these two constraints unless an explicit geometry or algebraic mechanism forces the identification.

\paragraph*{Section status.}  This section is a boundary condition for the broader project.  It prevents overclaiming about topology and generations.
